% Backup of §3.3 "Edits Given as Writes" (round-2 items 1-4 + round-1 rest), removed 2026-11-08 by decision
% (docs/s33-decision.md, option D). Its proof (Lemma 1) follows the subsection.
\subsection{Edits Given as Writes}
\label{sec:fw:edits}

% Spine: docs/framework-design.md, Section 13 (approved 2026-11-05 with the recommendations Q1-Q4). Revised in the
% review sprint (docs/sprint-framework-review.md): round 1 2026-11-08; round 2 in progress, items 1-4 (lead, written
% set, Algorithm 2 with structure positions and margins, the three cases, local kinds) drafted for the user's read.
An edit stated in items, such as one that adds an item with a given content, comes with its delta. Many edits are given
instead by the elements they write: they overwrite some elements, append new ones, or insert them.
\Cref{sec:fw:directions} needs to know which items such edits can change, for every content an attacker might write,
without trying any of them. This subsection shows that a write that leaves alone the positions an extractor reads to
locate items changes items only near the positions it writes. The delta of one write can then be read from the elements
near the write, and the items that any write within a region can change are known from the items of $z$ alone.

\definition[Written set]
\label{def:fw:written}
{
The \emph{written set} $W$ of an edit that takes $z$ to $z'$ is the set of positions at which $z$ and $z'$ differ,
including the positions that exist in only one of them.
}

\noindent Overwriting and appending have small written sets. Inserting an element in the middle moves every element
after it, so its written set runs to the end of the input.

\begin{algorithm}[t]
  \caption{How the extractor of \Cref{alg:fw:extractor} finds its records and tokens. A header at the start of $z$
  locates the records, and the tokens are found by scanning.}
  \label{alg:fw:finding}
  \begin{algorithmic}[1]
    \State $n \gets z[1]$ \Comment{header: the number of records}
    \For{$i = 1, \dots, n$}
      \State $o \gets z[2i]$, $\ell \gets z[2i + 1]$ \Comment{entry $i$: offset and length}
      \State record $i$ has the members $\{o, \dots, o + \ell - 1\}$
    \EndFor
    \State the tokens are the maximal runs of letters in $z$ \Comment{scanning}
  \end{algorithmic}
\end{algorithm}

\Cref{alg:fw:finding} shows the two ways in which an extractor finds its items. Records are \emph{located}: the
extractor reads a count and, for each record, an offset and a length, so the positions of the header decide where every
record lies. We call the positions that decide where the items of a kind $\tau$ lie its \emph{structure positions}
$L_\tau(z)$. For records they are the header, and tokens and elements have none. A write that changes no structure
position of records leaves every record where it was, and changes at most the content of the records it writes. Tokens
are \emph{scanned}: a write can lengthen a token, split it, or merge it with its neighbor, but only near the write,
since a run of letters cannot extend past an unchanged element that is not a letter. The tokens that can change
therefore lie between the nearest such element before the write and the nearest one after it. We call the region
around a write in which the items of $\tau$ can change the \emph{margin} $\mu_\tau(z, W)$. It depends on the input and
on the written set, it contains the written set and every item that has a member there, and a larger written set has a
larger margin. The margin of records is the written set together with the members of the records it writes, and that
of elements is the written set itself.

Every edit then falls into one of three cases. An edit stated in items comes with its delta. A write that changes no
structure position changes items only within the margins, and Lemma~\ref{lem:fw:localization} below reads its delta
there. A write that changes a structure position, such as the count of records or an offset, may change the items of
that kind anywhere. The framework does not predict what such a write does, and the end of this subsection says how it
is treated.

\definition[Local kinds]
\label{def:fw:local}
{
A kind $\tau$ is \emph{local} if it has structure positions $L_\tau(z)$ and margins $\mu_\tau(z, W)$ such that every
edit of $z$ whose written set $W$ contains no structure position of $\tau$ satisfies two conditions, where an item
\emph{lies within} a set of positions if all its members do:
\begin{enumerate}[label=(\roman*), leftmargin=*, nosep]
  \item the items of kind $\tau$ that do not lie within $\mu_\tau(z, W)$ are the same in $z$ and in $z'$;
  \item the items of kind $\tau$ of $z'$ that lie within $\mu_\tau(z, W)$ are determined by the elements of $z'$ on
    $\mu_\tau(z, W)$.
\end{enumerate}
Its margins must contain $W$ and the members of every item of kind $\tau$ that has a member in $W$, and must grow with
$W$.
}

% ROUND 2: the text below is still the round-1 version (items 5-8 of the round-2 plan come after the user's read).
\noindent Condition (ii) says that reading the margin of $z'$ again finds the items of the kind there, and the usual
ways of finding items meet both conditions with small margins. A kind found through fields, such as a record that a
count and an offset locate, has as structure positions the fields read to find its items, and its margin is $W$ with
the members of the items that $W$ meets. A write that misses these fields leaves every item where it was, and we say
that such a kind has a \emph{fixed layout}: under every write that misses its structure positions, its items in $z'$
have the same members as its items in $z$. A kind found by scanning the input for maximal runs of elements of some
class needs no structure positions, and its margin extends $W$ on each side to the nearest element of $z$ that lies
outside both the class and $W$. That element is unchanged in $z'$, so no run crosses it. For windows of $n$
consecutive elements, the margin extends $W$ by $n - 1$ elements on each side. These margins are the regions that an
incremental parser reads again after an edit, since it re-parses only the part of a program near the
change~\cite{wagner1998efficient}. An item whose members run to the end of the input, such as everything after the last
field, changes its members whenever a write reaches past the old end, so its margin then contains the whole item.

A write that misses the structure positions of a kind thus splits the items of that kind in two: those outside the
margin, which the edit keeps, and those inside, which the elements of $z'$ on the margin decide.

\lemma[Localization]
\label{lem:fw:localization}
{
Let $\tau$ be a local kind, and let an edit take $z$ to $z'$ with a written set $W$ that misses $L_\tau(z)$. The items
of kind $\tau$ of $z'$ are those of $z$, except that the items that lie within $\mu_\tau(z, W)$ are replaced by the
items of kind $\tau$ that the elements of $z'$ on the margin determine. Consequently:
\begin{enumerate}[label=(\roman*), leftmargin=*, nosep]
  \item if every kind is local and $W$ misses the structure positions of every kind, the kinds and contents of the
    replaced items and of their replacements form a delta of $(z, z')$, which follows, with the members of the
    replacements, from $z$, its item table, $W$, and the written elements;
  \item if a family of edits of $z$ writes only within a set $\hat{W}$ that misses $L_\tau(z)$, every edit of the
    family removes and adds only items of kind $\tau$ that lie within $\mu_\tau(z, \hat{W})$;
  \item if $\tau$ has a fixed layout, the edit only replaces items of kind $\tau$, each by an item with the same
    members.
\end{enumerate}
}

\noindent The main claim is Definition~\ref{def:fw:local} applied to one edit. Its first condition keeps the items
outside the margin, and their contents are unchanged as well: the margin contains the members of every item that $W$
meets, so the members of these items miss $W$. Its second condition says that the items inside are those the elements of
$z'$ on the margin determine. For part (i), the margins depend only on $z$ and $W$, the item table of $z$ lists the
items within them with their contents, and the elements of $z'$ on a margin are those of $z$ with the written elements
on $W$. Reading these elements finds the replacements with their members and, through the contribution functions, their
contributions, without building $z'$, and doing so for every kind gives the two multisets of a delta. Part (ii) holds
because every written set of the family lies in $\hat{W}$ and the margin grows with the written set. Part (iii) holds
because a fixed layout keeps the members of every item, so each replaced item has a replacement with the same members.
The input format enters only through the structure positions and the margins, which the instrumentation provides for
each kind, so applying the lemma to an extractor comes down to meeting Definition~\ref{def:fw:local}.

For a single write, part (i) gives the delta that Proposition~\ref{prop:fw:side-effects} needs, and also the members of
the new items, which a delta does not record. The item table of $z'$ is therefore known as well, and a sequence of
writes is followed by updating the item table one edit at a time. Part (ii) needs neither the content of the writes nor
any edited input. It bounds, from the item table of $z$ alone, which items of each kind a whole family of edits can
remove or add, and part (iii) says that a kind with a fixed layout can only be rewritten in place.
\Cref{sec:fw:directions} turns these bounds into directions.

A write that meets the structure positions of a kind can change the items of that kind anywhere, and the lemma says
nothing about them, although it still applies to the other kinds. Such an edit is handled when it is also stated in
items, so that it comes with its delta. A later write then also needs the members of the items the first edit added, so
an edit stated in items states their members along with their contents. A write that comes with no delta is
\emph{flagged}, and the items of the kinds whose structure positions it writes may change arbitrarily. Whether a write
is flagged is decided by comparing $W$ with the structure positions, which the instrumentation records together with the
item table. Definition~\ref{def:fw:local} excludes no kind. Every kind is local if all positions are its structure
positions, and a kind whose items cannot be found again from a small margin takes more positions as structure positions.
A large set of structure positions costs flags, since more writes are flagged, but never a wrong delta.


\paragraph{Proof of Lemma~\ref{lem:fw:localization}.}
Write $\mu = \mu_\tau(z, W)$, let $A$ be the items of kind $\tau$ of $z$ that lie within $\mu$, and $A'$ those of $z'$.
By condition (i) of Definition~\ref{def:fw:local}, the other items of kind $\tau$ are the same in $z$ and $z'$. Their
members miss $W$, since $\mu$ contains the members of every item of kind $\tau$ that $W$ meets, and $z'$ equals $z$
outside $W$, so their contents are the same too. By condition (ii), $A'$ is determined by the elements of $z'$ at the
positions of $\mu$. Hence the items of kind $\tau$ of $z'$ are those of $z$ with $A$ replaced by $A'$.
(i) If every kind is local and $W$ misses the structure positions of every kind, apply the main claim to each kind and
take the union: with $I^-$ and $I^+$ the multisets of the kinds and contents of the replaced items and of their
replacements, $I^- \subseteq \pi(z)$ and $\pi(z') = (\pi(z) - I^-) \uplus I^+$, which is \eqref{eq:fw:delta}. The
margins are functions of $z$ and $W$, and the replaced items, with their contents, and the members and contributions
of the other items are in the item table of $z$. The elements of $z'$ on the margins are those of $z$ outside $W$ and
the written elements on $W$, so the replacements, their members, and, through the contribution functions, their
contributions follow without building $z'$.
(ii) Every written set $W$ of the family lies in $\hat{W}$, so it misses $L_\tau(z)$ and, since the margin grows with
the written set, $\mu_\tau(z, W) \subseteq \mu_\tau(z, \hat{W})$. By the main claim, the edit removes only items that
lie within $\mu_\tau(z, W)$ and adds only items that lie there.
(iii) With a fixed layout, the items of kind $\tau$ of $z'$ have the same members as those of $z$. The items outside
$\mu$ are kept, so the members of the items in $A'$ are those of the items in $A$, and each replaced item is replaced
by the item of $z'$ with the same members.

